\section{Conclusion and Future Scope}
\label{ch:concl}
\providecommand{\HopW}{99.2}\providecommand{\HopH}{57.4}\providecommand{\HopWins}{20}\providecommand{\HopRuns}{20}

\providecommand{\SosAdv}{4.4}

\providecommand{\CollN}{50}\providecommand{\CollDel}{92.6}\providecommand{\CollDelMin}{76}\providecommand{\CollDelMax}{100}\providecommand{\CollRecMean}{8.1}\providecommand{\CollRecMax}{10.25}\providecommand{\CollRecN}{50}\providecommand{\CollRecOK}{23}\providecommand{\CollConvMean}{9.6}\providecommand{\CollConvMax}{15.2}


\subsection{Conclusion}
Every mainstream messaging tool fails when its infrastructure fails, although devices in a stricken building can still reach each other. We built EcoMesh, a peer-to-peer mesh in which every device is an endpoint and a relay, with a complete engine, a simulator that measures it, a REST/WebSocket interface, two front ends and a multi-device deployment over TCP. Weighted routing delivered \HopW\,\% of messages where hop-count routing delivered \HopH\,\% on identical seeds; an SOS had a \SosAdv$\times$ lower median latency under load; a 40-node mesh at 20\,\% loss with injected failures delivered \CollDel\,\% over 50 runs; and on five live processes an SOS arrived in a median of 32\,ms, the sender re-routed around a failed relay in about 4.1\,s, and every message held during a gateway outage was delivered when the gateway returned.

\subsection{Achievement of Objectives}
Objectives O1--O3 and O5--O8 were achieved. O4 was achieved for delivery but only partly for recovery time, where 26 of 50 trials exceeded 10\,s by at most 0.25\,s, the heartbeat timeout (Table~\ref{tab:objectives}).

\subsection{Major Engineering Contributions}
A transport-agnostic engine that runs unchanged on a simulator and on real sockets; a derived link-quality cost function with an immediate failure penalty and network-wide re-flooding; per-hop priority with a bounded queue that never evicts an SOS for routine traffic; signatures that survive multi-hop relaying; and a measurement framework in which every number in this report is generated from per-node logs and stored results.


\subsection{Current Limitations}
As detailed in Section~\ref{sec:limits}: a neighbour-restricted LAN stands in for a radio; one pre-shared key serves the mesh; relays do not retain; timer constants quantise recovery and convergence; there is no usability study with participants, and the browser checks are scripts run by hand rather than part of the automated suite; there is no phone application, only a web page that a phone opens over the same Wi-Fi; and scale was tested to 65 simulated nodes.


\subsection{Future Enhancements}
(1) A Bluetooth LE or Wi-Fi Direct transport and an Android client speaking the same protocol. (2) Adaptive timers: a shorter fixed neighbour timeout (about 7\,s) was tried and dropped because it lowered delivery at 20\,\% loss, so the next step is a timeout adapted to each link, an immediate report on topology change and an ACK timeout derived from the measured round-trip time. (3) X25519 per-pair key agreement. (4) Short retention at relays, so that a relay does not lose a message it has accepted. (5) Per-destination route invalidation. (6) A usability study, running the browser checks inside the automated suite, a stable ordering for messages created in the same millisecond (D7), and tests beyond 65 nodes.

\subsection{Lessons Learned}
Build the simulator first and draw the transport boundary before the layers: it paid for itself in testing and in deployment. Integration finds what unit tests do not (the multi-hop ACK bug survived every single-node test). Look at distributions as well as means: recovery clustered at the heartbeat timeout, which the mean hid. Record who built what from the first week, with consistent Git identities.

